\chapter{WAV vs MP3 και Codec Matching}
\label{chap:codec-wav-mp3}

Πριν από κάθε εκπαίδευση, το σύστημα πρέπει να διασφαλίσει ότι
το format των αρχείων ήχου δεν αποτελεί shortcut για τον ταξινομητή.
Το παρόν κεφάλαιο τεκμηριώνει το πρόβλημα (100\% συσχέτιση format–label),
τη λύση (MP3 round-trip του bonafide μέσω \codeid{src/data/codec\_match\_audio.py}),
και την επέκταση της Phase~6+ για καθαρούς-WAV spoof providers (\engineName{OpenAI~TTS}).

% ============================================================
\section{Το Πρόβλημα: 100\% Format-Label Correlation}
\label{sec:format-label-correlation}
% ============================================================

Τα αρχεία ήχου που εισέρχονται στο pipeline δεν έχουν ομοιόμορφο format.
Κάθε κλάση εγγραφής αποθηκεύεται σε διαφορετικό codec
ανάλογα με την πηγή της:

\begin{table}[ht]
  \centering
  \begin{tabular}{lll}
    \toprule
    Row class & On disk & Why \\
    \midrule
    bonafide real & clean 16 kHz PCM WAV & από \texttt{ffmpeg} extract \\
    ElevenLabs spoof & lossy MP3 44.1k/128k & TTS engine output \\
    Google spoof & lossy MP3 24k/64k & TTS engine output \\
    OpenAI spoof & clean WAV & default \texttt{response\_format=wav} \\
    \bottomrule
  \end{tabular}
  \caption{Pre-codec-match format ανά row class.}
  \label{tab:precodec}
\end{table}

Η κατανομή του Πίνακα~\ref{tab:precodec} αποκαλύπτει ένα κρίσιμο πρόβλημα:
ο format του αρχείου ήχου είναι \textbf{100\% συσχετισμένος με την ετικέτα}.
Κάθε μοντέλο που δέχεται mel-spectrogram ή raw waveform ως είσοδο
μπορεί θεωρητικά να μάθει να ανιχνεύει τη διαφορά PCM~WAV vs.\ MP3
αντί για πραγματική deepfake ανίχνευση.
Μια τέτοια στρατηγική θα έδινε υψηλή ακρίβεια στο training και validation set,
αλλά θα κατέρρεε σε οποιοδήποτε dataset όπου το TTS output
είναι clean WAV — ή αντίστροφα, όπου τα bonafide αρχεία
έχουν περάσει από MP3 επεξεργασία.

Αυτό το φαινόμενο ονομάζεται \textbf{codec leak}:
ένα τεχνητό αρχείο-format signal διαρρέει στο training loop
ως ψεύτικο discriminative cue.
Ακόμα και μετά τη χρήση frozen Wav2Vec2/WavLM/HuBERT embeddings
αντί για raw mel-spectrograms, το codec history εξακολουθεί
να αφήνει ίχνη στη φασματική κατανομή των hidden states.

% ============================================================
\section{Λύση: MP3 Round-Trip του Bonafide}
\label{sec:mp3-roundtrip}
% ============================================================

Η λύση είναι να εξομοιωθεί το codec history των bonafide εγγραφών
με αυτό των spoof εγγραφών.
Αντί να αφαιρεθεί ο MP3 codec από τα spoof αρχεία
(αδύνατο — η lossy συμπίεση είναι μη αντιστρέψιμη),
τα bonafide αρχεία \textbf{round-tripάρονται μέσω MP3}:
κωδικοποιούνται σε MP3 με τις παραμέτρους ενός native-MP3 spoof provider,
και αμέσως αποκωδικοποιούνται πίσω σε 16~kHz mono PCM WAV.
Έτσι κάθε bonafide αρχείο αποκτά κι αυτό codec history,
και το format του αρχείου παύει να είναι predictive του label.

% ------------------------------------------------------------
\subsection{Deterministic Codec Sampling}
\label{subsec:deterministic-codec-sampling}
% ------------------------------------------------------------

Για κάθε bonafide εγγραφή, ο codec spec που θα χρησιμοποιηθεί
για το MP3 round-trip επιλέγεται \textbf{ντετερμινιστικά} —
όχι stochastically — μέσω της συνάρτησης
\codeid{codec\_for\_bonafide(sample\_id, provider\_weights)}.

Ο αλγόριθμος είναι απλός και αναπαραγώγιμος:

\begin{english}
\begin{enumerate}
  \item Υπολόγισε τον SHA-1 hash του \codeid{sample\_id} ως ακέραιο
        (\codeid{h = int(hashlib.sha1(sample\_id.encode()).hexdigest(), 16)}).
  \item Πάρε το υπόλοιπο της διαίρεσης με το συνολικό πλήθος
        native-MP3 spoof εγγραφών:
        \codeid{point = h \% total}.
  \item Διάτρεξε τα native-MP3 providers σε ταξινομημένη σειρά
        (lexicographic, για αναπαραγωγιμότητα),
        συσσωρεύοντας τα πλήθη τους.
        Επίστρεψε τον provider για τον οποίο
        \codeid{point < cum} (cumulative count).
\end{enumerate}
\end{english}

Το αποτέλεσμα είναι ότι η marginal κατανομή του codec spec
πάνω από όλες τις bonafide εγγραφές
\textbf{συγκλίνει στην κατανομή του native-MP3 spoof corpus},
ανεξάρτητα από τη σειρά επεξεργασίας των εγγραφών.
Ο SHA-1 modulo cumulative distribution εξασφαλίζει
αναπαραγωγιμότητα μεταξύ εκτελέσεων —
η ανάθεση codec σε κάθε \codeid{sample\_id} δεν αλλάζει ποτέ.

Τα specs αυτά προέκυψαν από \codeid{ffprobe} ανάλυση των on-disk αρχείων
και είναι σταθερά ανά provider.

% ------------------------------------------------------------
\subsection{Implementation Details}
\label{subsec:implementation-details}
% ------------------------------------------------------------

Η υλοποίηση βρίσκεται στο \codeid{src/data/codec\_match\_audio.py}.
Η κεντρική συνάρτηση \codeid{process\_row()} χειρίζεται κάθε εγγραφή ως εξής:

\begin{english}
\begin{enumerate}
  \item Αν η εγγραφή είναι bonafide ή ανήκει σε \codeid{CLEAN\_WAV\_PROVIDERS}:
    \begin{enumerate}
      \item Κάλεσε \codeid{codec\_for\_bonafide()} για να αποκτήσεις
            \codeid{(provider, sr, bitrate)}.
      \item Κωδικοποίησε το αρχείο WAV σε MP3 με
            \codeid{ffmpeg -c:a libmp3lame -ar sr -b:a bitrate}
            σε ένα temporary directory.
      \item Αποκωδικοποίησε το MP3 πίσω σε 16~kHz mono PCM WAV
            (\codeid{ffmpeg -ar 16000 -ac 1 -c:a pcm\_s16le}).
      \item Διαγράφει το temporary MP3.
    \end{enumerate}
  \item Αν η εγγραφή είναι native-MP3 spoof:
    \begin{enumerate}
      \item Αποκωδικοποίησε απευθείας το MP3 σε 16~kHz mono PCM WAV,
            διατηρώντας το υπάρχον codec history.
    \end{enumerate}
\end{enumerate}
\end{english}

Η έξοδος αποθηκεύεται στον κατάλογο
\codeid{data/audio\_wav\_codec\_matched/\{sample\_id\}.wav}.
Παράλληλα, γράφεται ένα derived manifest
(\codeid{data/derived/audio\_spoof\_manifest\_codec\_matched.csv})
όπου η στήλη \codeid{audio\_path} δείχνει στα νέα WAV αρχεία,
ώστε τα downstream embedding extractors να μπορούν να τα καταναλώσουν
χωρίς τροποποίηση.

Ο χειρισμός σφαλμάτων γίνεται per-row:
αν το \codeid{ffmpeg} αποτύχει για μια εγγραφή,
η εγγραφή καταγράφεται ως \codeid{failed} και η επεξεργασία συνεχίζεται.
Αυτό επιτρέπει partial runs χωρίς να διακόπτεται ολόκληρο το pipeline.

% ============================================================
\section{Phase 6+ Extension: Clean-WAV Spoof Providers}
\label{sec:clean-wav-extension}
% ============================================================

% ------------------------------------------------------------
\subsection{Γιατί το OpenAI Δημιουργεί Νέο Confounder}
\label{subsec:openai-confounder}
% ------------------------------------------------------------

Η Phase~6+ εισήγαγε τον \engineName{OpenAI~TTS} ως τρίτο TTS engine.
Αντίθετα με τα \engineName{ElevenLabs} και \engineName{Google~TTS}
που παράγουν MP3 by default,
το \engineName{OpenAI~TTS} χρησιμοποιεί
\codeid{response\_format=wav} ως προεπιλογή —
δηλαδή παράγει clean PCM WAV χωρίς MP3 codec history.

Αν εφαρμοζόταν το codec-match \emph{μόνο} στα bonafide αρχεία
(όπως στην αρχική Phase~4 λύση),
το \engineName{OpenAI~TTS} θα παρέμενε ο μοναδικός clean-WAV spoof provider.
Αυτό θα δημιουργούσε έναν νέο codec leak —
\textbf{αντίστροφης κατεύθυνσης}:
αντί το μοντέλο να μαθαίνει «clean WAV → bonafide (real)»,
θα μπορούσε να μάθει «clean WAV → OpenAI (spoof)».

Η φόρμουλα του προβλήματος αλλάζει σημείο,
αλλά η δομή παραμένει ίδια:
ένα αρχείο-format signal διαρρέει ως discriminative cue
αντί για γνήσια deepfake ανίχνευση.
Η διόρθωση πρέπει να εφαρμοστεί και στις δύο clean-WAV κλάσεις.

% ------------------------------------------------------------
\subsection{\codeid{CLEAN\_WAV\_PROVIDERS = \{"openai\_tts"\}}}
\label{subsec:clean-wav-providers}
% ------------------------------------------------------------

Κάθε spoof εγγραφή της οποίας ο provider ανήκει στο
\codeid{CLEAN\_WAV\_PROVIDERS} υποβάλλεται στο \textbf{ίδιο MP3 round-trip}
με τις bonafide εγγραφές.
Κρίσιμη λεπτομέρεια: το sampling distribution για το codec spec
\textbf{περιορίζεται αποκλειστικά στα native-MP3 providers}
(εκείνους που ανήκουν στο \codeid{PROVIDER\_CODEC}).
Αυτό εξασφαλίζει ότι το marginal codec footprint
για τα bonafide και τα OpenAI rows
ταυτίζεται με την κατανομή των native-MP3 providers
(\engineName{ElevenLabs} + \engineName{Google~TTS}).
Το αποτέλεσμα: ο codec δεν αποτελεί πλέον predictive signal
\emph{για καμία} κλάση — ούτε για τα bonafide,
ούτε για τους native-MP3 spoofs,
ούτε για τους clean-WAV spoofs.

% ============================================================
\section{Επαλήθευση στο Πραγματικό Run}
\label{sec:codec-verification}
% ============================================================

Η εκτέλεση του \codeid{codec\_match\_audio} στο Phase~6+ manifest
(4\,149 bonafide + 2\,218 spoof, συνολικά 6\,367 εγγραφές)
Τα νούμερα αυτά επιβεβαιώνουν τη λογική:

\begin{itemize}
  \item \textbf{5\,196 εγγραφές round-tripάρουν μέσω MP3}:
        οι 4\,149 bonafide και οι 1\,047 OpenAI spoof εγγραφές.
  \item \textbf{2\,218 εγγραφές} (native-MP3 spoof)
        αποκωδικοποιούνται απευθείας σε 16~kHz WAV χωρίς νέο encoding.
  \item Το codec spec sampling βαρύνεται
        673 ElevenLabs vs.\ 498 Google:
        περίπου 57\% των round-trip εγγραφών
        λαμβάνουν spec \codeid{(44100, 128k)}
        και 43\% \codeid{(24000, 64k)}.
  \item Ο συνολικός αριθμός εξόδων (\codeid{written})
        ισούται με 6\,367 — κάθε εγγραφή του manifest
        παράγει ακριβώς ένα \codeid{.wav} αρχείο
        στον κατάλογο \codeid{data/audio\_wav\_codec\_matched/}.
\end{itemize}

Το round-trip είναι ένα \textbf{μοναδικό ντετερμινιστικό βήμα},
όχι stochastic augmentation.
Ο SHA-1 modulo cumulative distribution εξασφαλίζει
ότι η ανάθεση codec spec σε κάθε \codeid{sample\_id}
είναι ταυτόσημη σε κάθε εκτέλεση —
ανεξάρτητα από τη σειρά εκτέλεσης των rows,
την έκδοση του Python, ή άλλους περιβαλλοντικούς παράγοντες.
Αυτό καθιστά το codec-match βήμα
πλήρως αναπαραγώγιμο ως μέρος του frozen pipeline.
